\section{Related Work}
\label{sec:related}

\paragraph{Kinematic signatures of driver impairment.}
The link between vehicle kinematics and driver cognitive state has a firm theoretical and empirical basis.
Steering entropy theory, introduced by Nakayama et al.~\cite{Nakayama1999} and formalised by Boer~\cite{Boer2000},
establishes that degraded closed-loop vehicle control under cognitive load manifests as increased irregularity
in steering corrections.
Empirical support is extensive: Strayer and Johnston showed that even hands-free phone conversation
impairs driving relative to an undistracted baseline~\cite{Strayer2001};
Strayer et al.\ further demonstrated that cell-phone use produces speed and following-distance
degradation comparable to legal alcohol impairment~\cite{Strayer2003,Strayer2006};
Recarte and Nunes confirmed that mental workload reduces the attentional bandwidth available for
vehicle control independently of visual distraction~\cite{Recarte2003};
and Patten et al.\ showed that even hands-free mobile phone use imposes measurable cognitive
costs on primary driving performance~\cite{Patten2004}.
Engstr\"{o}m et al.\ showed that visual and cognitive distraction produce distinct, identifiable
signatures across speed regulation, lane-keeping precision, and lateral acceleration~\cite{Engstrom2005};
Kountouriotis et al.\ confirmed that steering wheel reversal rates are selectively sensitive
to cognitive secondary task load even when gaze remains on the road~\cite{Kountouriotis2016};
and Klauer et al.\ demonstrated, in the landmark 100-Car Naturalistic Driving Study, that
inattention is the primary contributing factor in near-crash events~\cite{Klauer2006}.
A comprehensive theoretical account is provided by Engstr\"{o}m et al.~\cite{Engstrom2017},
who establish that cognitive load systematically reduces the bandwidth of the driver's
closed-loop vehicle control, producing observable kinematic consequences before overt errors occur.
More recently, Zheng et al.\ showed that increases in steering entropy predict the onset of performance
degradation in a driving simulator~\cite{Zheng2022}.
The same kinematic sensitivity extends to driver fatigue: Li et al.\ demonstrated online drowsiness
detection from steering wheel angles under real driving conditions~\cite{LiFatigue2017},
and Arefnezhad et al.\ achieved accurate drowsiness classification using adaptive neuro-fuzzy
feature selection applied exclusively to steering data~\cite{Arefnezhad2019}.
Taken together, this body of evidence motivates the use of kinematic signals as a proxy
for impairment state---encompassing both cognitive distraction and fatigue---rather than
as predictors of discrete error events.

\paragraph{Signal-based driver distraction detection.}
Leveraging these kinematic signatures for automatic detection has attracted sustained interest.
Li et al.\ demonstrated that driving performance indicators derived from steering, speed, and
lateral position in a large-scale naturalistic dataset can discriminate distracted from
attentive driving with meaningful accuracy, and highlighted steering entropy as the most
informative single feature~\cite{Li2018}.
Kaplan et al.\ provide a comprehensive survey of signal-based driver behaviour analysis,
noting that threshold-based and handcrafted-feature approaches achieve reasonable within-subject
performance but fail to generalise across drivers due to large inter-individual variability
in baseline kinematics~\cite{Kaplan2015}.
Dingus et al.\ analysed the SHRP2 naturalistic driving study and confirmed that
driver-related factors are implicated in approximately 88\% of crashes,
reinforcing the value of continuous driver state monitoring---while also demonstrating
the rarity and context-dependence of discrete crash events, which limits the
viability of event-prediction framing~\cite{Dingus2016}.

\paragraph{Machine learning for driver state monitoring.}
Classical machine learning approaches---Support Vector Machines, Random Forests,
and Gradient Boosted Trees---have been widely applied to driver state classification
from vehicle signals, typically operating on hand-engineered time-domain and
frequency-domain features~\cite{Kaplan2015, Kumar2023, Xue2023}.
Xue et al.\ directly compared SVM and Random Forest on vehicle control signals
for cognitive distraction detection, finding that both achieve competitive
accuracy within-subject but neither was evaluated across
participants~\cite{Xue2023}.
Deep learning methods, particularly Long Short-Term Memory
networks~\cite{Hochreiter1997} and CNN-LSTM hybrids~\cite{Kwon2021}, have
emerged as the dominant paradigm for end-to-end learning from raw or lightly
preprocessed driving time series.
Wang et al.\ demonstrated that a Bi-LSTM with attention trained on vehicle
dynamics alone achieves strong distraction detection performance on naturalistic
driving data~\cite{Wang2022}, reinforcing the viability of kinematic-only deep
learning approaches while again relying on within-subject evaluation.
Temporal Convolutional Networks (TCNs), introduced as a competitive
sequence modelling alternative to recurrent architectures by Bai et al.~\cite{Bai2018},
combine causal dilated convolutions~\cite{vandenOord2016} with residual
connections~\cite{He2016} to achieve receptive fields that scale exponentially
with depth, enabling efficient capture of long-range temporal dependencies.
Despite their demonstrated advantages in generic sequence modelling
tasks~\cite{Bai2018}, TCNs remain underexplored in the driver monitoring
literature relative to LSTM-based approaches~\cite{Kaplan2015, Wang2022}.
Zhao et al.\ applied a soft-thresholding TCN to abnormal driving behaviour
recognition from vehicle dynamics, reporting improved efficiency and
generalisation over recurrent baselines~\cite{Zhao2022},
but the evaluation was conducted within-subject.
To our knowledge, no prior work has applied a TCN to driver impairment scoring
from vehicle kinematic signals alone under a fully cross-participant evaluation protocol.

\paragraph{Evaluation rigour and the cross-participant generalisation gap.}
A persistent methodological concern in the driver monitoring literature is the
reliance on within-subject evaluation protocols, in which training and test windows
are drawn from the same participants~\cite{Kaplan2015}.
This practice systematically inflates reported performance, because models can
exploit individual-specific driving style rather than learning generalisable
impairment signatures.
Kaplan et al.\ note this limitation explicitly, observing that kinematic-based
systems that appear effective in within-subject settings do not transfer reliably
to new drivers~\cite{Kaplan2015}.
The problem is compounded when studies also perform data augmentation or
normalisation using information from the test participant---a form of leakage
that further overstates generalisation.
Leave-One-Participant-Out Cross-Validation (LOPO-CV) is the strictest available
protocol for addressing this gap: each participant's data is held out entirely,
and all preprocessing and model fitting is performed exclusively on the remaining
participants~\cite{Engstrom2017}.
Our work adopts LOPO-CV as the sole evaluation protocol, applies anti-leakage
route renormalisation to prevent any test-participant statistics from influencing
the training fold, and additionally reports false positive rates on a held-out
cohort of participants with no induced errors---a safety-oriented audit that has
not been reported in prior kinematic-based driver monitoring studies.
